\subsection*{V6. Origin of Ferromagnetic Domains|Need, Energy terms|Formation \& Shape, Equations with Explanation}
\addcontentsline{toc}{subsection}{6. Origin of Ferromagnetic Domains|Need, Energy terms|Formation \& Shape, Equations with Explanation}

\textcolor{red}{Q.No. 1 Why does iron pieces do not attract each other? why do they need a magnet to attract each other?}\\
\ding{228} This can be explained by the concept of domains.
\begin{enumerate}
    \item Weiss postulated that the ferromagnetic substance is divided into a large number of small domains. In each of these domains, the magnetization is aligned in one direction. The direction of magnetization in each domain are opposite and tends to cancel each other.
    \item Each dimension is magnetized but the directions of magnetism in various domains tends to cancel each other.
    \item The Net magnetization is zero.
    \item \textcolor{red}{\textbf{How to see Domains?}}
    \begin{enumerate}
        \item It was developed by Francis Bitter in 1931.
        \item In this method, the surface of the sample is covered with an aqueous solution of very small colloidal particles of magnetite $\mathrm{(Fe_3O_4)}$.
        \item Magnetite deposits as a band along the domain boundaries, at their intersection with the sample surface.
        \item Outlines of domains can be seen using a microscope.
        \item In general we can see the domains by following steps:
        \begin{enumerate}
            \item Carefully polish the surface of ferromagnet.
            \item Spread a fine powder of ferromagnetic particles over it.
            \item Particles collect over domain boundaries.
            \item Observed using different microscopy methods. Some of them are listed below.
        \end{enumerate}
    \end{enumerate}
    \item \textbf{Visualization of Magnetic Domains:}
    \begin{enumerate}
        \item Images and details of different microscopy like, Kerr Microscope, Magnetic Force Microscopy, and Lorentz Transmission electron microscopy are just said and images are provided in this section.
    \end{enumerate}
    \item \textbf{Formation \& shape of Domains:}
    \begin{enumerate}
        \item The formation and the shape of domains depends on the competition among different energy terms present in the magnetic crystal.
        \item Consider a uniformly magnetized state of a crystal, which means the exchange energy is lowest, and it can be said that this have large amount of magneto-static energy. 
        \item Magneto-static energy is a self energy due to interaction of magnetic field due to magnetization in some part of the sample on the other parts of the same sample.
        \item This energy depends on the volume occupied by magnetic field extending outside the domain.
        \item As in the uniformly magnetized charge which is on the surface of the sample. These charges are going to produce the magnetic field which will be in the opposite direction to that of the magnetization and that field which opposes the magnetization is called the demagnetization field which will be always opposite to $\mathrm{\overrightarrow{M}}$ denoted as $\mathrm{\overrightarrow{B_d}}$.
        \item Hence, we can say that the net magnetization would be the competition between these two fields.
        \item $\mathrm{\overrightarrow{M}}$ is opposite to $\mathrm{\overrightarrow{B_d}}$.
        \item Magneto-static energy is positive.
        \item The corresponding density is given as,
        \begin{align}
            E_m &= \frac{1}{2} M B_d 
            \intertext{where, the demagnetizing field is given by,}
            \vec{B}_d &= -\mu_0 D \vec{M}
        \end{align}
         where, $D$ is the demagnetization factor and is large for flat sample and small for elongated sample.
        and demagnetization energy is given as
        \item We can see that the demagnetization field is itself dependent on the the magnetization and it comes into the picture only when there is surface magnetic charges present in the system.
        \item 
    \end{enumerate}
    \item  
\end{enumerate}